%% ============================================================================
%% main_zh.tex, Chinese mirror of main_en.tex.
%% Build: xelatex main_zh && bibtex main_zh && xelatex main_zh && xelatex main_zh
%% Mirror discipline: same sections, same \label names, same theorem numbering, same citation keys and the
%% same \num keys as main_en.tex. check_mirror.py verifies this.
%% ============================================================================
\documentclass[10pt]{article}
\usepackage[letterpaper,margin=1in]{geometry}

%% ---- Chinese typesetting ---------------------------------------------------
%% Font fallback chain. FandolSong ships with TeX Live, so the build is reproducible on any machine. The next
%% three are the usual system serifs on macOS, Windows and Linux.
\usepackage{xeCJK}
\IfFontExistsTF{FandolSong-Regular.otf}{%
  \setCJKmainfont{FandolSong-Regular.otf}[BoldFont=FandolSong-Bold.otf,ItalicFont=FandolKai-Regular.otf]%
  \setCJKsansfont{FandolHei-Regular.otf}[BoldFont=FandolHei-Bold.otf]%
  \setCJKmonofont{FandolFang-Regular.otf}%
}{%
  \IfFontExistsTF{Songti SC}{%
    \setCJKmainfont{Songti SC}[BoldFont=Heiti SC,ItalicFont=Kaiti SC]%
  }{%
    \IfFontExistsTF{SimSun}{%
      \setCJKmainfont{SimSun}[BoldFont=SimHei,ItalicFont=KaiTi]%
    }{%
      \setCJKmainfont{Noto Serif CJK SC}%
    }%
  }%
}
\XeTeXlinebreaklocale "zh"
\XeTeXlinebreakskip = 0pt plus 1pt

%% ---- Chinese head names, declared before the shared preamble ---------------
\newcommand{\headtheorem}{定理}
\newcommand{\headlemma}{引理}
\newcommand{\headproposition}{命题}
\newcommand{\headcorollary}{推论}
\newcommand{\headdefinition}{定义}
\newcommand{\headassumption}{假设}
\newcommand{\headexample}{例}
\newcommand{\headremark}{注}
\newcommand{\headfinding}{发现}
\newcommand{\headproof}{证明}

%% ============================================================================
%% preamble_shared.tex
%% Shared preamble of main_en.tex and main_zh.tex. Everything that must be identical across the two language
%% versions lives here, so that section files, labels, theorem counters, number macros and notation stay
%% mirrored by construction. main_zh.tex defines the Chinese head names before it reads this file.
%% ============================================================================

%% ---- Latin text font -------------------------------------------------------
%% Under pdflatex, main_en.tex loads the times package. Under xelatex, TeX Gyre Termes supplies the Times
%% metrics when it is installed. fontspec is never loaded under pdflatex.
\usepackage{iftex}
\ifPDFTeX
  \usepackage[T1]{fontenc}
\else
  \usepackage{fontspec}
  \IfFontExistsTF{texgyretermes-regular.otf}{%
    \setmainfont{texgyretermes}[Extension=.otf, UprightFont=*-regular, BoldFont=*-bold,
      ItalicFont=*-italic, BoldItalicFont=*-bolditalic]%
  }{}
\fi

%% ---- packages --------------------------------------------------------------
\usepackage{amsmath,amssymb,amsthm}
\usepackage{booktabs,array,adjustbox,multirow,colortbl,pifont}
\usepackage{graphicx}
\usepackage{xcolor}
\usepackage{algorithm}
\usepackage{algorithmic}
\renewcommand{\algorithmicrequire}{\textbf{Input:}}
\renewcommand{\algorithmicensure}{\textbf{Output:}}
\usepackage{subcaption}
\usepackage{url}
\usepackage{hyperref}
\definecolor{citecolor}{HTML}{009B3A}
\definecolor{refcolor}{HTML}{4527C4}
\hypersetup{colorlinks=true, citecolor=citecolor, linkcolor=refcolor, urlcolor=refcolor, pdfborder={0 0 0}}

%% ---- table style (canonical copy in figure-style/tables/table_macros.tex) --
% Table style of figure-style/tables, one palette shared by the main, capability and dataset tables.
% Put % Table style of figure-style/tables, one palette shared by the main, capability and dataset tables.
% Put \input{table_macros} in the preamble. Needs booktabs, colortbl, xcolor, array, pifont.
%
%   first place in a column   \best{0.957}         dark background, bold
%   second place in a column  \secondbest{0.847}   light background
%   family row of main table  \grp{\ncol}{Deep matching} \\
%   row of the proposed method \oursrow before the row
%   capability marks          \cmark  \xmark
%   value waiting for backfill \pend{0.xx}          red, must be gone before submission

\definecolor{tabBest}{RGB}{31,78,121}     % first place, used at 35 percent tint
\definecolor{tabSecond}{RGB}{207,224,237} % second place
\definecolor{tabBand}{RGB}{226,234,242}   % row of the proposed method

\newcommand{\best}[1]{\cellcolor{tabBest!35}\textbf{#1}}
\newcommand{\secondbest}[1]{\cellcolor{tabSecond}#1}
\newcommand{\grp}[2]{\multicolumn{#1}{l}{\textbf{\textit{#2}}}}
\newcommand{\oursrow}{\rowcolor{tabBand}}
\newcommand{\cmark}{\ding{51}}
\newcommand{\xmark}{\ding{55}}
\newcommand{\pend}[1]{\textcolor{red}{#1}}
 in the preamble. Needs booktabs, colortbl, xcolor, array, pifont.
%
%   first place in a column   \best{0.957}         dark background, bold
%   second place in a column  \secondbest{0.847}   light background
%   family row of main table  \grp{\ncol}{Deep matching} \\
%   row of the proposed method \oursrow before the row
%   capability marks          \cmark  \xmark
%   value waiting for backfill \pend{0.xx}          red, must be gone before submission

\definecolor{tabBest}{RGB}{31,78,121}     % first place, used at 35 percent tint
\definecolor{tabSecond}{RGB}{207,224,237} % second place
\definecolor{tabBand}{RGB}{226,234,242}   % row of the proposed method

\newcommand{\best}[1]{\cellcolor{tabBest!35}\textbf{#1}}
\newcommand{\secondbest}[1]{\cellcolor{tabSecond}#1}
\newcommand{\grp}[2]{\multicolumn{#1}{l}{\textbf{\textit{#2}}}}
\newcommand{\oursrow}{\rowcolor{tabBand}}
\newcommand{\cmark}{\ding{51}}
\newcommand{\xmark}{\ding{55}}
\newcommand{\pend}[1]{\textcolor{red}{#1}}

\makeatletter\let\inputrows\@@input\makeatother

%% ---- theorem-like environments ---------------------------------------------
%% Head names go through macros, so both language versions share counters and numbering.
\providecommand{\headtheorem}{Theorem}
\providecommand{\headlemma}{Lemma}
\providecommand{\headproposition}{Proposition}
\providecommand{\headcorollary}{Corollary}
\providecommand{\headdefinition}{Definition}
\providecommand{\headassumption}{Assumption}
\providecommand{\headexample}{Example}
\providecommand{\headremark}{Remark}
\providecommand{\headfinding}{Finding}
\providecommand{\headproof}{Proof}

\theoremstyle{plain}
\newtheorem{theorem}{\headtheorem}[section]
\newtheorem{lemma}{\headlemma}[section]
\newtheorem{proposition}{\headproposition}[section]
\newtheorem{corollary}{\headcorollary}[section]
\theoremstyle{definition}
\newtheorem{definition}{\headdefinition}[section]
\newtheorem{assumption}{\headassumption}[section]
\newtheorem{example}{\headexample}[section]
%% Findings are numbered through the whole paper (Finding 1 to Finding 7), not per section.
\newtheorem{finding}{\headfinding}
\theoremstyle{remark}
\newtheorem{remark}{\headremark}[section]
\renewcommand{\proofname}{\headproof}

%% ---- system name -----------------------------------------------------------
\newcommand{\sys}{\mbox{OurSys}}

%% ---- measured numbers --------------------------------------------------------
%% Every measured number is \num{key}. templates/backfill/backfill.py writes the values into numbers_<lang>.tex
%% from the result files, so a rerun refills both papers. A key without a value prints in red.
\newcommand{\pending}[1]{\textcolor{red}{[\detokenize{#1}]}}
\newcommand{\setnum}[2]{\expandafter\gdef\csname rbnum@#1\endcsname{#2}}
\newcommand{\num}[1]{\ifcsname rbnum@#1\endcsname\csname rbnum@#1\endcsname\else\pending{#1}\fi}

%% ---- notation (docs/CONCEPTS.md, one concept one name one symbol) ----------
\newcommand{\budget}{B}
\newcommand{\cand}{\mathcal{C}}
\newcommand{\policy}{\pi}

% Generated by backfill.py from registry.example.json. Do not edit by hand (zh).
\setnum{baselines}{4}
\setnum{datasets}{3}
\setnum{families}{2}
\setnum{gain over best pct}{5.9}
\setnum{label share pct}{5}
\setnum{recall best baseline}{0.861}
\setnum{recall ours}{0.912}
% waiting: runtime minutes <- demo_results/timing.json largest_table_minutes (file missing)


\renewcommand{\abstractname}{摘要}
\renewcommand{\refname}{参考文献}
\renewcommand{\figurename}{图}
\renewcommand{\tablename}{表}
\renewcommand{\appendixname}{附录}
\floatname{algorithm}{算法}

\title{\sys{}: <一句点明主要结果的话>}
\author{匿名作者}
\date{}

\begin{document}
\maketitle

\begin{abstract}
% 六到八句，每句对应 knowledge/paper-anatomy.md 第 1 节摘要表的一行。最强的事实放在最后一两句。
<对象>在<场景>中承担<作用>，而<某种性质>带来了问题。
<朴素做法>不能说明<期望结果>，所以<决策>应当由<正确的判据>来评判。
本文提出 \sys{}，它<步骤 1>、<步骤 2>，再<步骤 3>。
我们证明，在<条件>下，<保证>(定理~\ref{thm:guarantee})。
在 \num{datasets} 个数据集上，\sys{} 只用 \num{label share pct}\% 的标注就达到 \num{recall ours} 的召回率，\num{baselines} 个基线中最强者为 \num{recall best baseline}。
\end{abstract}

\section{引言}
\label{sec:intro}
% 五步，见 knowledge/paper-anatomy.md 第 2 节。不设小标题，用 \paragraph{}。

% (1) 代价。一个真实事故或调查数字，再收成开放问题。
<问题>普遍存在且代价高昂~\cite{placeholder2024}。<年份>，<带数字的具体事故>~\cite{placeholder2024}。如何在只有<约束>时<目标>，仍是一个开放问题。

% (2) 视角转换与命名的性质。
然而<传统目标>并不能保证<下游目标>。在这一视角下，两个性质共同决定<结果>。\emph{<性质 A>}度量<……>。\emph{<性质 B>}度量<……>。图~\ref{fig:motivation}与例~\ref{ex:motivation}具体展示了二者的取舍。

\begin{figure}[t]
\centering
\fbox{\parbox{0.9\linewidth}{\centering\vspace{1.2cm}动机图，见 knowledge/figure-archetypes.md 第 1 节。\vspace{1.2cm}}}
\caption{<在什么条件下比较了什么，每个面板一句自然的话>。}
\label{fig:motivation}
\end{figure}

% (3) 真实数据集上的例子与 Observation。
\begin{example}[<短标题>]
\label{ex:motivation}
我们以公开数据集<数据集>~\cite{placeholder2024}的一个子集为例。<每行是什么、任务和模型>。<划分与规模>。<搜索空间以及图~\ref{fig:motivation}展示了什么>。
\end{example}

\textbf{观察。}由这个例子可以得到两点观察。
\textbf{(1) <一句话论断>。}<带数字的证据>。
\textbf{(2) <一句话论断>。}<带数字的证据>。

% (4) 挑战。一个承上启下的问句，再列编号的挑战。
\paragraph{挑战。}上述观察说明<要求>。那么现有方法为何没有解决<问题>？原因可归结为三项挑战。
\textbf{(1) <收益依赖任务>。}<原因，结合例~\ref{ex:motivation}>。现有的<方法族>~\cite{placeholder2024}缺少<能力>。
\textbf{(2) <决策空间组合爆炸>。}<原因>。
\textbf{(3) <决策难以迁移>。}<原因>。

% (5) 贡献。第一条承载核心论点，最后一条讲实验。
\paragraph{贡献。}针对上述不足，本文提出 \sys{}，<一句话定义>，贡献如下。
(1) \textbf{<名称>。}<做了什么、为什么有效，定理~\ref{thm:guarantee}>。
(2) \textbf{<名称>。}<做了什么、为什么有效，第~\ref{sec:method}~节>。
(3) \textbf{实验结果。}我们在 \num{datasets} 个数据集上与 \num{baselines} 个基线比较。\sys{} <头条结果>，代价为<代价>。

\section{相关工作}
\label{sec:related}
% 每个方法族一段，段首加粗族名，顺序与表~\ref{tab:related}和主表一致。

\textbf{基于规则的方法(RB)。}<代表方法>~\cite{placeholder2024}<做了什么>。它们<不能做什么>。

\textbf{基于学习的方法(LB)。}<代表方法>~\cite{placeholder2024}<做了什么>。它们<不能做什么>。

\begin{table}[t]
\centering
\caption{各对比方法利用什么、输出什么、提供什么保证。}
\label{tab:related}
\footnotesize
\setlength{\tabcolsep}{2.2pt}
\begin{tabular}{@{}>{\raggedleft\arraybackslash}p{0.34\linewidth}|c|cccc@{}}
\toprule
Method & Fam. & Bud. & Typ. & Task & Guar. \\
\midrule
<Method A>~\cite{placeholder2024} & RB & \xmark & \cmark & \xmark & \xmark \\
<Method B>~\cite{placeholder2024} & LB & \cmark & \xmark & \xmark & \xmark \\
\midrule
\oursrow
\textbf{\sys} & & \cmark & \cmark & \cmark & \cmark \\
\bottomrule
\end{tabular}
\end{table}

\section{问题定义}
\label{sec:problem}
% 输入对象、决策、策略、优化问题。每个定义后跟贯穿例子。

\begin{definition}[候选集]
\label{def:candidate}
给定表 $T$，候选集 $\cand$ 是<定义>。
\end{definition}

\begin{example}
\label{ex:candidate}
在例~\ref{ex:motivation}中，<贯穿例子的候选集>。
\end{example}

\begin{definition}[标注策略]
\label{def:policy}
策略 $\policy$ 把<状态>映射为<决策>，且至多标注 $\budget$ 个候选。
\end{definition}

目标为
\begin{equation}
\label{eq:objective}
\max_{\policy}\ \mathrm{Quality}(\policy(\cand)) \quad \text{s.t.} \quad |\policy(\cand)| \le \budget .
\end{equation}

\begin{theorem}[困难性]
\label{thm:hard}
在一般的逐候选代价下，问题~\eqref{eq:objective}的判定版本是 NP 完全的。
\end{theorem}
证明见附录~\ref{app:proof-hard}。这一结果说明第~\ref{sec:method}~节需要近似方法。

\section{\sys{} 方法}
\label{sec:method}
% 跟着系统图走。每一步依次写目的、定义、算法、保证、复杂度。
\sys{} 由两步组成。第 1 步<……>。第 2 步<……>，其输出反馈给下一轮的第 1 步。

\subsection{第 1 步：<名称>}
\label{sec:step1}
<一句话说明目的>。<定义>。<在贯穿例子上的演示>。

\subsection{第 2 步：<名称>}
\label{sec:step2}
<一句话说明目的>。算法~\ref{alg:main}列出了具体过程。

\begin{algorithm}[t]
\caption{<过程名称>}
\label{alg:main}
\begin{algorithmic}[1]
\REQUIRE 候选集 $\cand$，预算 $\budget$
\ENSURE 已标注集合 $L$
\STATE $L \gets \emptyset$
\WHILE{$|L| < \budget$}
  \STATE <按第 2 步的规则选择下一个候选>
\ENDWHILE
\RETURN $L$
\end{algorithmic}
\end{algorithm}

\begin{theorem}[保证]
\label{thm:guarantee}
在假设 <k> 下，算法~\ref{alg:main}的策略不低于<它所包含的每个基线>。
\end{theorem}
% 实现不满足定理的全部前提时，把定理写成条件保证。

\section{实验评估}
\label{sec:experiments}
我们回答三个研究问题。RQ 1 考察 \sys{} 是否<……>(贡献 1)。RQ 2 考察<……>。RQ 3 考察<……>。实验覆盖 \num{families} 个方法族、\num{baselines} 个基线和 \num{datasets} 个数据集。

\subsection{实验设置}
\label{sec:setup}
(1)\textbf{数据集。}表~\ref{tab:datasets}列出了各数据集的任务、错误类型与规模。<数据从哪来、错误从哪来>。
(2)\textbf{基线。}我们与表~\ref{tab:related}中的全部方法比较。<所有方法从什么数据学习>。<\sys{} 有何不同>。
(3)\textbf{指标。}我们报告\textbf{召回率}(\textbf{R})，$R = \tfrac{|\text{找到的重复}|}{|\text{全部重复}|}$，以及<……>。

\begin{table}[t]
\centering
\caption{各数据集的任务、错误类型与规模。}
\label{tab:datasets}
\footnotesize
\begin{tabular}{@{}r|c|c|c@{}}
\toprule
Dataset & Task & Error types & Records \\
\midrule
\textbf{<D1>}~\cite{placeholder2024} & DD & T/M & <n> \\
\textbf{<D2>}~\cite{placeholder2024} & DD & T/F & <n> \\
\bottomrule
\end{tabular}
\end{table}

\subsection{实验结果与发现}
\label{sec:findings}

\begin{table}[t]
\centering
\caption{相同标注预算下各方法的召回率。}
\label{tab:main}
\footnotesize
\setlength{\tabcolsep}{3.6pt}
\renewcommand{\arraystretch}{0.9}
\begin{tabular}{@{}r@{\hspace{4pt}}r|ccc@{}}
\toprule
Rank & Method & D1$\uparrow$ & D2$\uparrow$ & D3$\uparrow$ \\
\midrule
%%BEGIN-MAINROWS
\grp{5}{Rule-based} \\
4 & R-Match & 0.612 & 0.574 & 0.655 \\
5 & R-Block & 0.598 & 0.603 & 0.621 \\
\grp{5}{Learning-based} \\
3 & L-Forest & 0.844 & 0.792 & \secondbest{0.861} \\
2 & L-Boost & \secondbest{0.851} & \secondbest{0.805} & 0.833 \\
\midrule
\oursrow
1 & \textbf{\sys} & \best{0.912} & \best{0.866} & \best{0.890} \\
%%END-MAINROWS
\bottomrule
\end{tabular}
\end{table}

\begin{finding}
\label{find:main}
\sys{} <带有它所胜出的比较的强论断>。
\end{finding}
\sys{} 比最强的基线高 \num{gain over best pct}\%(表~\ref{tab:main})。<机理，对应引言的观察 (1)>。<边界情形及其实测原因>。

\begin{finding}
\label{find:cost}
<关于代价或规模的论断>。
\end{finding}
\sys{} 在最大的表上用时 \num{runtime minutes} 分钟。<哪一步占主要开销，随规模如何增长>。

\section{结论}
\label{sec:conclusion}
% 从实验事实起手，写本文证明了什么，再写适用条件和下一步。不复述数字，不重复摘要。
<\sys{} 说明了什么，一两句>。<适用条件>。<下一步>。


\bibliographystyle{plain}
\bibliography{refs}

\appendix
\section{定理~\ref{thm:hard}的证明}
\label{app:proof-hard}
% 先用一句话复述结论，再证明。名称与正文逐字一致。
\begin{proof}
<从一个已知的 NP 完全问题归约>。
\end{proof}

\section{参数}
\label{app:params}
% 正文不写的全部取值，包括超参数、划分、种子和硬件。
\begin{table}[h]
\centering
\caption{全部实验的参数。}
\label{tab:params}
\footnotesize
\begin{tabular}{@{}ll@{}}
\toprule
参数 & 取值 \\
\midrule
<名称> & <取值> \\
\bottomrule
\end{tabular}
\end{table}

\end{document}
